% Source PDF 15–17, printed 3–5; all three scans visually checked.
\chapter*{引言}\addcontentsline{toc}{chapter}{引言}
% SOURCE-PAGE-ANCHOR: 15
\sourcepage{15}
下降理论在代数几何学家中多少有些令人望而生畏的名声. 事实上, 它只是说, 在一定条件下, 拟凝聚层之间的同态可以先局部构造, 再在满足相容条件时黏合起来; 而拟凝聚层本身也可以先局部构造, 再通过满足余循环条件的同构黏合起来.

当然, 如果“局部”指“在 Zariski 拓扑下局部”, 这只是一个形式性的陈述, 固然有用, 却很难称得上一个理论. 要点在于, 这里的“局部”指在平坦拓扑下局部; 而平坦拓扑并不是通常意义的拓扑, 而是所谓的 \emph{Grothendieck 拓扑}. 在这里, 覆盖本质上是满足某个有限性条件的平坦满射. 因而这种拓扑中的覆盖远多于 Zariski 拓扑中的覆盖, 证明也就变得很不平凡.

不过, 只要采用通常的滥用记号, 把层 $F$ 沿两个映射 $f$ 与 $g$ 的复合所得的拉回 $(gf)^*F$ 与 $f^*g^*F$ 认同, 陈述仍然十分简单自然. 若要完全严格, 就必须考虑到 $(gf)^*F$ 与 $f^*g^*F$ 并不相同, 而只是有一个满足某些相容条件的典范同构, 并且必须发展这些相容性的理论. 在我看来, 由此产生的复杂性, 正是许多代数几何学家厌弃下降理论的原因（如果他们还愿意看它一眼的话）.

还有一个抽象概念, 即“下降理论成立的范畴”; 由概形与其上拟凝聚层组成的二元组所成范畴就是一个例子. 这些范畴称为\emph{叠}（stacks）. 这一一般形式体系即使在模理论之外也相当有用; 在模理论中, 代数叠理论已处于绝对核心地位（例如参见 [\textbf{DM69}]、[\textbf{Art74b}] 和 [\textbf{LMB00}]）.

这些讲义原是为我在 2003 年 7 月 7 至 18 日于 I.C.T.P. 举办的 \textit{Advanced School in Basic Algebraic Geometry} 中关于 \textit{Grothendieck topologies and descent theory} 的十次讲课所写. 它们旨在比 Grothendieck 的原始论文（[\textbf{Gro95b}], 至今仍很好读, 强烈推荐）或 [\textbf{SGA1}] 更完整地阐述下降理论. 我也采用了在这里很自然的 Grothendieck 拓扑语言, 但在上述两种标准文献写成时, 这一语言还没有被引入.

这里的论述侧重纤维范畴与叠的一般理论: 希望立刻得到应用的代数几何学家可能会感到失望. 不过, 我认为一般理论既有趣又适用, 也希望至少有一部分读者会赞同.

% Source PDF 16, printed 4.
% SOURCE-PAGE-ANCHOR: 16
\sourcepage{16}
此外, 对拟凝聚层和概形的下降理论的讨论, 是这些讲义真正的存在理由. 在其中, 我从不采用“只要有典范同构就将对象认同”的约定, 而总是明确指出同构, 并写出所需的相容条件. 这样使论述严格, 但也相当繁重（一个格外令人不快的例子见 §4.3.3）. 人们或许会怀疑这种选择是否明智; 但我想说服自己, 完全严格的处理确实是可能的. 那位不幸的读者可以放心: 这给我带来的痛苦比给她带来的更多.

这些讲义中的所有思想与结果都归功于 Grothendieck. 这里没有什么东西不是以某种形式出现在 [\textbf{SGA1}] 或 [\textbf{SGA4}] 中的, 因而我完全不声称有任何原创性.

下降理论还有一些现代发展, 尤其是在范畴论中（例如参见 [\textbf{JT84}]）, 以及非交换代数与非交换几何中（[\textbf{KR04a}] 和 [\textbf{KR04b}]）. 对拓扑学家和代数几何学家都最激动人心的方向之一, 是与高阶范畴论这一重要主题紧密相连的“高阶下降”思想（例如参见 [\textbf{HS}] 和 [\textbf{Str03}]）. 我们不会讨论这些非常有趣的主题.

\medskip\noindent\textbf{内容.}
第 1 章回顾代数几何与范畴论的一些基本概念.

真正的内容从第 2 章开始. 我先讨论 Grothendieck 的可表函子观点, 并给出一个主要的说明性例子: 展示这一观点如何使群概形以及群概形在概形上的作用变得十分自然、容易. 全部代数几何都可以从这一观点系统发展, 从而变得非常整洁优美, 而又令初学者难以理解（参见 [\textbf{DG70}]）.

在第 2.3 节, 我定义并讨论 Grothendieck 拓扑及其上的层. 我采用预拓扑的朴素观点, 觉得它直观得多. 不过, 使用筛的较精细观点也有优点, 所以我试图两全其美（意大利语的说法更生动:“既让我的酒桶满满, 又让我的妻子喝醉”）: 定义筛, 并用筛刻画层, 从而隐含地说明层只依赖拓扑, 不依赖预拓扑. 本节还介绍概形范畴上的四种主要拓扑: Zariski、平展、fppf 和 fpqc, 并证明 Grothendieck 的定理: 可表函子在这些拓扑下都是层.

下降理论有两种可能的形式框架: 纤维范畴和伪函子. 前者似乎较不繁琐, 因此第 3 章讨论纤维范畴理论. 不过, 我在这里也定义伪函子, 并联系这两种观点, 因为若干例子（如拟凝聚层）用伪函子的语言表达更加自然. 我证明一些重要结果（首先是纤维范畴的 Yoneda 引理）, 最后讨论纤维范畴中的等变对象（希望有些读者会觉得, 这有助于理解相当复杂的等变层概念）.

这些讲义的核心是第 4 章. 在充分讨论下降数据之后（我给出几种定义, 并证明其等价）, 定义核心概念——\emph{叠}: 叠是一个带有 Grothendieck 拓扑的范畴上的纤维范畴, 在其中下降理论成立（这样我们便看到了标题中三个概念如何共同发挥作用）. 随后证明主定理: 拟凝聚层的纤维范畴在 fpqc 拓扑下是叠. 然后将其用于概形的下降理论成立的两个主要例子, 即仿射态射以及配有典范丰沛线丛的态射. 我还讨论一个特别有趣的例子, 即沿主丛的下降（用 Grothendieck 的术语说, 挠子 torsors）.

% Previous paragraph spans source PDF 16–17; remaining source PDF 17.
% SOURCE-PAGE-ANCHOR: 17
\sourcepage{17}
最后一节给出一个例子, 说明平展下降对概形并不总是成立; 最后提及概形概念的一种推广, 即 Michael Artin 提出的\emph{代数空间}. 它的用处在于: 一方面, 代数空间在某种意义上非常接近概形, 可以为其定义概形论中的大多数概念; 另一方面, fppf 下降对代数空间总是成立. 本来将它收入讲义也是很自然的, 但那样会进一步推迟讲义的完成.

\medskip\noindent\textbf{先修要求.}
我假定读者熟悉概形的语言, 至少达到 Hartshorne 的书 [\textbf{Har77}] 的水平. 我会用到其中没有的一些概念, 例如局部有限呈示态射; 但第 1 章会回顾其主要性质, 并给出 \textit{Éléments de géométrie algébrique} 中相应部分的文献出处.

我大量采用范畴语言: 假定读者熟悉范畴、函子、自然变换和范畴等价的概念. 另一方面, 我不使用高深概念, 也不使用范畴论中实质性的结果, 只有一个例外: 读者应当知道, 全忠实且本质满射的函子是等价.

\medskip\noindent\textbf{致谢.}
在 \textit{Advanced School in Basic Algebraic Geometry} 讲授这门课程是一段非常愉快的经历, 这要归功于其他讲课者（Lothar Göttsche、Luc Illusie、Steve Kleiman 和 Nitin Nitsure）的友谊, 以及学员们积极热情的态度. 我还感谢 Lothar、Luc、Steve 和 Nitin, 因为他们从未抱怨过这些讲义迟迟未能完成.

感谢 Steve Kleiman 提供有益的讨论和建议, 尤其是关于 fpqc 拓扑的部分; 也感谢 Pino Rosolini, 在几次阿尔卑斯山徒步途中, 他尝试为我讲明某些范畴构造.

我与 Behrang Noohi 就叠的定义有过一些有趣的交谈: 在此向他表示衷心感谢.

我从 Andrew Kresch 那里得知了 [\textbf{Ray70}, XII 3.2] 中的反例.

我也感谢许多对我的系列讲课表示兴趣的学员, 特别是指出初版讲义错误的那些人. 尤其感谢 Zoran Skoda, 他寄来了若干有益的意见, 还在参考文献方面给予帮助.

Joachim Kock 仔细阅读了其中大部分内容, 寄来一长列非常有用的意见与订正. 审稿人又提供了更多订正. 我感谢他们.

最后, 我愿将这些讲义献给我的岳父 Amleto Rosolini, 以纪念他. 讲义完成之际, 他去世了, 享年 86 岁. 他本来不会有兴趣学习下降理论, 但他为人善良而出众, 对数学的热情一直持续到生命的最后一天, 这将永远激励我.
